\section{Modifications to chapter 6}

The following modifications shall be applied to chapter 6 (\emph{Basics}
[\texttt{basic}]):

\begin{modification}[Modification of 6.1/5]
Every name is introduced by a \emph{declaration}, which is a

\begin{itemize}[label=\textemdash]
\item \ldots

\item
\stdremove{\emph{identifier} in a \emph{result-name-introducer} in a
postcondition assertion (9.4.2),}

\item \ldots
\end{itemize}

\end{modification}

\begin{modification}[Modification of 6.1/8]
An entity is a variable, structured binding, \stdremove{result binding,} function, 
enumerator, type, type alias, non-static
data member, bit-field, template, namespace, namespace alias, template parameter, function parameter, or
init-capture. The underlying entity of an entity is that entity unless otherwise specified. A name denotes the
underlying entity of the entity declared by each declaration that introduces the name.
\end{modification}

\begin{modification}[Modification of 6.1/9]
A \emph{local entity} is a variable with automatic storage duration (6.8.6.4), a structured binding (9.7) whose
corresponding variable is such an entity, \stdremove{a result binding (9.4.2),} or the *this object (7.5.3).
\end{modification}

\begin{modification}[Modification of 6.3/10.2]
\begin{itemize}[label=\textemdash]

\item
for each intervening scope (6.4.1) between the point at which the entity is introduced and the scope
(where \texttt{*this} is considered to be introduced within the innermost enclosing class or non-lambda function
parameter scope), either

\begin{itemize}[label=\textemdash]

\item
the intervening scope is a block scope,

\item
\stdremove{the intervening scope is a contract-assertion scope (6.4.10),}

\item
the intervening scope is the function parameter scope of a lambda-expression or requires-expression,
or

\item
the intervening scope is the lambda scope of a lambda-expression that has a simple-capture naming
the entity or has a capture-default, and the block scope of the lambda-expression is also an intervening
scope.
\end{itemize}
\end{itemize}
\end{modification}

\begin{modification}[Modification of 6.4.1/1]
The declarations in a program appear in a number of scopes that are in general discontiguous. The global scope
contains the entire program; every other scope \emph{S} is introduced by a
declaration, \emph{parameter-declaration-clause},
\emph{statement}, \emph{handler}, \stdadd{or}
\emph{lambda-expression}\stdremove{, or \emph{contract assertion}} (as described in the following subclauses of 6.4)
appearing in another scope, which thereby contains \emph{S}. An enclosing scope at a program point is any scope
that contains it; the smallest such scope is said to be the immediate scope at that point. A scope intervenes
between a program point \emph{P} and a scope \emph{S} (that does not contain
\emph{P}) if it is or contains \emph{S} but does not
contain \emph{P}.
\end{modification}

\begin{modification}[Modification of 6.4.1/5]
A declaration is \emph{name-independent} if its name is \texttt{\_} (U+005F LOW
LINE) and it declares

\begin{itemize}[label=\textemdash]

\item
a variable, other than a function parameter, with automatic storage duration,

\item
a structured binding with no \emph{storage-class-specifier} and not inhabiting a namespace scope,

\item
\stdremove{a result binding (9.4.2),}

\item
the variable introduced by an \emph{init-capture}, or

\item
a non-static data member of other than an anonymous union.

\end{itemize}

\emph{Recommended practice}: Implementations should not emit a warning that a name-independent declaration is
used or unused.
\end{modification}

\begin{modification}[Removal of 6.4.2/13]
\stdremove{The locus of a result-name-introducer (9.4.2) is immediately after it.}
\end{modification}

\begin{modification}[Removal of 6.4.10]
Section 6.4.10 (\emph{Contract-assertion scope} [\texttt{basic.scope.contract}])
shall be \textbf{completely removed}.
\end{modification}

\begin{modification}[Modification of 6.8.6.5.2/5]
A global allocation function is only called as the result of a new expression (7.6.2.8), or called directly using
the function call syntax (7.6.1.3), or called indirectly to allocate storage for a coroutine state (9.6.4), or called
indirectly through calls to the functions in the C++ standard library.
\newline

NOTE 2 In particular, a global allocation function is not called to allocate storage for objects with static storage
duration (6.8.6.2), for objects or references with thread storage duration
(6.8.6.3), for objects of type \texttt{std::typeinfo} (7.6.1.8)
\stdremove{, for an object of type 
\texttt{std::contracts::contract\_violation} when a contract violation occurs
(6.11.2),}
or for an exception object (14.2).
\end{modification}

\begin{modification}[Modification of 6.10.1/5]
A \emph{full-expression} is

\begin{itemize}[label=\textemdash]
\item
an unevaluated operand (7.2.3),

\item
a constant-expression (7.7.2),

\item
an immediate invocation (7.7.5),

\item
an init-declarator (9.3) (including such introduced by a structured binding (9.7)) or a mem-initializer
(11.9.3), including the constituent expressions of the initializer,

\item
an invocation of a destructor generated at the end of the lifetime of an object other than a temporary
object (6.8.7) whose lifetime has not been extended, \stdadd{or}

\item
\stdremove{the predicate of a contract assertion (6.11), or}

\item
an expression that is not a subexpression of another expression and that is not otherwise part of a
full-expression.
\end{itemize}

\end{modification}

\begin{modification}[Modification of 6.10.1/11]
When invoking a function \emph{f} (whether or not the function is inline), every argument expression and the postfix
expression designating \emph{f} are sequenced before \stdremove{every
precondition assertion of \emph{f} (9.4.1), which in turn are
sequenced before} every expression or statement in the body of f\stdremove{, which in turn are sequenced before every
postcondition assertion of \emph{f}}.
\end{modification}

\begin{modification}[Modification of 6.10.3.1/5]
A return statement (8.8.4) in \texttt{main} has the effect of leaving the main function (destroying any objects with
automatic storage duration \stdremove{and evaluating any postcondition
assertions of \texttt{main}}) and calling \texttt{std::exit} with
the return value as the argument. If control flows off the end of the compound-statement of main, the effect
is equivalent to a return with operand \texttt{0} (see also 14.4).
\end{modification}

\begin{modification}[Removal of 6.11]
Section 6.11 (\emph{Contract assertions} [basic.contract]) shall be \textbf{completely
removed}.
\end{modification}
